\clearpage
\VerdierPhysicalPageAnchor{0205}%
\begingroup\small\setlength{\parskip}{0pt}\setlength{\emergencystretch}{1.5em}
\setlength{\abovedisplayskip}{4pt plus 2pt minus 2pt}%
\setlength{\belowdisplayskip}{4pt plus 2pt minus 2pt}%

\VerdierSemanticAnchor{ch3d-2-4-4-proof-a-conclusion}
which shows that $X$ is a projective object of $\mathcal A$.

\VerdierSemanticAnchor{ch3d-2-4-4-proof-b}
Let us prove b). Let $L^\bullet$ be a complex of projective type that is an
object of $K^{-}(\mathcal A)$, and let
$u\colon L^\bullet\longrightarrow Y$ be a morphism into an acyclic
complex. Suppose that, for every integer $i>n$, we have constructed a
morphism $s^i\colon L^i\longrightarrow Y^{i-1}$ such that
$u^i=d_Y^{i-1}s^i+s^{i+1}d_{L^\bullet}^{i}$ (such a construction is always
possible for $n$ sufficiently large because
$L^\bullet$ is an object of $K^{-}(\mathcal A)$). It follows that
$d_Y^n\bigl(u^n-s^{n+1}d_{L^\bullet}^{n}\bigr)=0$, and therefore the
morphism $u^n-s^{n+1}d_{L^\bullet}^{n}$ factors through
$\operatorname{Ker}(d_Y^n)$. Since the complex $Y$ is acyclic, the
morphism $Y^{n-1}\longrightarrow\operatorname{Ker}(d_Y^n)$ is an
epimorphism; and since $L^n$ is projective, there is a morphism
$s^n\colon L^n\longrightarrow Y^{n-1}$ such that
$d_Y^{n-1}s^n=u^n-s^{n+1}d_{L^\bullet}^{n}$.
Induction thus shows that, for every integer $i$, there is a morphism $s^i$
such that $u^i=d_Y^{i-1}s^i+s^{i+1}d_{L^\bullet}^{i}$. The morphism $u$
is therefore homotopic to zero. It follows from (Chap. II, 2.3.3) that
$L^\bullet$ is an object that is $Q$-free on the left. The last assertion
of b) now follows from (Chap. II, 2.3.3).

\VerdierSemanticAnchor{ch3d-2-4-4-proof-c-projective-complexes}
Let us prove c). Assertion 1) follows immediately from b) and (2.3.1). To
prove assertion 1) in its ``respectively'' form, first let us show that
every complex of projective type is $Q$-free on the left. Let $L^\bullet$
be a complex of projective type and let
$u\colon X\longrightarrow L^\bullet$ be a quasi-isomorphism. It
suffices to show that $u$ has a section (Chap. II, 2.3.3). It
therefore suffices to show that there is a complex ${L'}^\bullet$ of
projective type and a quasi-isomorphism
$v\colon {L'}^\bullet\longrightarrow X$ such that the composite morphism
$uv\colon {L'}^\bullet\longrightarrow L^\bullet$ is an isomorphism
in $K(\mathcal A)$. By (2.3.3), there is always a complex ${L'}^\bullet$
of projective type and a quasi-isomorphism
$v\colon {L'}^\bullet\longrightarrow X$. It remains only to show that the
quasi-isomorphism
$uv\colon {L'}^\bullet\longrightarrow L^\bullet$ is an isomorphism.
Now there is a distinguished triangle
$({L'}^\bullet,L^\bullet,{L''}^\bullet,uv,w,m)$, where
${L''}^\bullet$ is an acyclic complex of projective type.

\VerdierSemanticAnchor{ch3d-2-4-4-proof-c-acyclic-projective-complexes}
It remains only to show that acyclic complexes of projective type are
homotopic to zero. Let ${L''}^\bullet$ be an acyclic complex of projective
type. Since the projective dimension of the objects of $\mathcal A$ is
bounded, the kernels of the differentials of ${L''}^\bullet$ are
projective (2.2.6). It follows that all the differentials of
${L''}^\bullet$ are decomposable morphisms (Chap. II, 1.2.8), and hence
that ${L''}^\bullet$ is homotopic to zero.

\VerdierSemanticAnchor{ch3d-2-4-4-proof-c-conclusion}
It now follows from (Chap. II, 2.3.3) that the natural functor
$K^{*}(P)\longrightarrow D(\mathcal A)$ is fully faithful. Consequently,
the functor $K^{*}(P)\longrightarrow D^{*}(\mathcal A)$ is fully faithful.
By (2.3.3), this functor is also essentially surjective, and is therefore
an equivalence of categories. Assertions 2) and 3) are now merely direct
applications of (Chap. II, 2.3.3). Assertion 4) is then evident: an object
that is $Q$-free on the left in $K^{-}(\mathcal A)$ (respectively, in
$K^{*}(\mathcal A)$) admits, by (2.3.1) (respectively, (2.3.3)), a left
resolution by a complex of projective type in $K^{-}(\mathcal A)$
(respectively, $K^{*}(\mathcal A)$), and a quasi-isomorphism between
objects that are $Q$-free on the left is an isomorphism in $K(\mathcal A)$.
\endgroup
